\documentclass[11pt,a4paper]{article}
\usepackage[utf8]{inputenc}
\usepackage[vietnamese,english]{babel}
\usepackage[top=2cm,bottom=2cm,left=2.2cm,right=2.2cm]{geometry}
\usepackage{xcolor}
\usepackage{booktabs}
\usepackage{hyperref}

\definecolor{hcmutblue}{RGB}{0,62,126}
\definecolor{darkgray}{RGB}{60,60,60}

\hypersetup{colorlinks=true, linkcolor=hcmutblue, urlcolor=hcmutblue}

\begin{document}
\selectlanguage{vietnamese}

\begin{center}
  {\small \textbf{ĐẠI HỌC QUỐC GIA TP. HỒ CHÍ MINH -- TRƯỜNG ĐẠI HỌC BÁCH KHOA}}\\[0.05cm]
  {\small \textbf{KHOA KHOA HỌC VÀ KỸ THUẬT MÁY TÍNH}}\\[0.1cm]
  {\large\bfseries\color{hcmutblue} BẢN TUYÊN BỐ SỬ DỤNG CÔNG CỤ TRÍ TUỆ NHÂN TẠO}\\[0.05cm]
  {\large\bfseries\color{hcmutblue} (AI-USE STATEMENT)}\\[0.2cm]
  {\normalsize \textbf{Môn học:} CO5151 -- Advanced Agentic AI (HK261)}\\
  {\normalsize \textbf{Chủ đề Seminar:} S1-7: Agent Design Patterns \& Orchestration}\\
  {\normalsize \textbf{Nhóm thực hiện (G4):} Đặng Lâm Tùng (Trưởng nhóm) $\cdot$ Nguyễn Trung Phong $\cdot$ Vũ Việt Hưng}\\
  {\normalsize \textbf{Giảng viên hướng dẫn:} TS. Lê Xuân Bách}
\end{center}

\vspace{0.4cm}
\noindent Nhóm G4 cam kết tuân thủ nghiêm ngặt các quy định về liêm chính học thuật của Trường Đại học Bách khoa -- ĐHQG-HCM và hướng dẫn của môn học CO5151. Dưới đây là bảng kê khai chi tiết, minh bạch việc sử dụng các công cụ trí tuệ nhân tạo (AI Tools) trong toàn bộ quá trình chuẩn bị tài liệu Seminar S1-7:

\vspace{0.3cm}
\begin{table}[h]
\centering
\small
\begin{tabular}{p{3.2cm}p{4.2cm}p{6.8cm}}
\toprule
\textbf{Công cụ AI} & \textbf{Khâu sử dụng} & \textbf{Mô tả chi tiết tác vụ} \\
\midrule
\textbf{Google Gemini 3.8 Flash} & Hỗ trợ rà soát ngôn ngữ \& định dạng & Rà soát ngữ pháp các câu nhận xét tiếng Anh trong Annotated Bibliography; chuẩn hóa cú pháp bảng biểu LaTeX Beamer. \\
\addlinespace
\textbf{Claude 3.5 Sonnet (Cursor IDE)} & Hỗ trợ lập trình khung mẫu (Boilerplate) & Tạo khung mã nguồn mẫu cho kịch bản mô phỏng LangGraph Execution Trace (\texttt{demo\_orchestration\_trace.py}). \\
\bottomrule
\end{tabular}
\end{table}

\vspace{0.3cm}
\noindent \textbf{\large Các khâu KHÔNG sử dụng AI (100\% thành viên nhóm G4 tự thực hiện):}
\begin{enumerate}
  \item \textbf{Xây dựng luận điểm phản biện cốt lõi:} Toàn bộ nội dung phân tích bản chất Workflows vs. Agents, lan truyền lỗi $P=p^n$, chi phí bùng nổ token $3\times-8\times$ của Lead/Subagent và cơ chế Checkpointing bền bỉ được hình thành trực tiếp từ việc nghiên cứu sâu 4 tài liệu gốc (Anthropic 2024, Anthropic 2025, OpenAI 2025, LangGraph) đối chiếu với thực tế triển khai đề tài Theme 9 (Legal RAG).
  \item \textbf{Viết nhận xét phản biện từng tài liệu:} Từng câu nhận xét phản biện trong bản \textit{Annotated Bibliography} được phân công cụ thể cho từng thành viên đọc hiểu sâu và tự đúc kết lỗ hổng phương pháp luận.
  \item \textbf{Thiết kế ma trận so sánh các design patterns:} Nhóm tự phân loại các mẫu thiết kế và xác định các chiều tiêu chí kỹ thuật so sánh dựa trên bản chất kiến trúc hệ thống.
  \item \textbf{Kiểm chứng sự kiện (Fact-Checking):} Toàn bộ các công thức toán học, trích dẫn tác giả, mốc thời gian và số liệu thực nghiệm đều được kiểm tra chéo với tài liệu gốc chính thức.
\end{enumerate}

\vspace{0.8cm}
\begin{flushright}
  \textit{TP. Hồ Chí Minh, ngày 18 tháng 10 năm 2026}\\[0.2cm]
  \textbf{Đại diện Nhóm G4 (Trưởng nhóm)}\\[1.5cm]
  \textbf{Đặng Lâm Tùng}
\end{flushright}

\end{document}
